\section{离散系统与z域分析}\label{sec:Discrete_System_z_Analysis}

%本节中，我们将介绍离散系统的基本特性，并介绍由z变换引出的离散系统的系统函数，以及由系统函数引出的离散系统的频率响应和相位响应等概念。

\subsection{离散系统概述}\label{sub:discrete_system_overview}

离散系统是指激励与响应均为离散信号的系统。本小节中，我们先列举一些离散系统的基本特性，然后讨论离散系统的级联。

\begin{definition}
    设离散系统的激励为$e[n]$，响应为$r[n]=H(e[n])$，则该系统具有以下基本特性：
    \begin{itemize}
        \item \textbf{线性性}：满足叠加定理，即对于任意激励$e_1[n],e_2[n]$及任意常数$a,b\in\mathbb{C}$，有
        \[H(ae_1[n]+be_2[n])=aH(e_1[n])+bH(e_2[n])\]
        \item \textbf{时不变性}：系统的作用不依赖于时间，即对于任意激励$e[n]$及任意整数$k\in\mathbb{Z}$，有
        \[H(\tau_k e[n])=\tau_k H(e[n])\]
        其中$\tau_k$表示时间平移算子，即
        $$\tau_k e[n]=e[n-k]$$
        \item \textbf{因果性}：系统的响应不早于激励，即
        \[\forall e[n],e[n]=0,n<0\implies r[n]=0,n<0\]
        \item \textbf{稳定性}：激励信号有界时，响应也是有界的，即
        \[\forall e[n]\in l^{\infty}(\mathbb{Z}),r[n]\in l^{\infty}(\mathbb{Z})\]
        \item \textbf{可逆性}：如果$H$是双射，则称系统可逆，此时存在一个系统$H^{-1}$使得对于任意激励$e[n]$，有
        \[H^{-1}r[n]=e[n]\]
        \item \textbf{即时性}与\textbf{记忆性}：如果该系统在任意时刻$t$的状态完全且唯一地取决于该时刻的系统状态与激励，而不显式地依赖于任何过去时刻$t'<t$的系统状态或历史激励，则称该系统为即时的；否则称该系统为有记忆的。
    \end{itemize}
    其中，同时满足线性性和时不变性的系统称为线性时不变系统；即时性又称无记忆性、Markov性；记忆性又称历史依赖性、非Markov性。
\end{definition}

显然，线性系统的级联仍然是线性的，时不变系统的级联仍然是时不变的，因果系统的级联仍然是因果的。对于离散系统，“对核积分”就相当于加权求和，“对核卷积”就相当于卷积和，可以想象，叠加定理对于离散系统仍成立。

\begin{definition}[离散线性系统的脉冲响应]
    离散线性系统$H$的\textbf{核函数}定义为
    $$k[m,n]=H_m(\delta[m-n])$$
    其中$H_m$表示$\delta[m-n]$作为$m$的函数输入系统，输出为$m,n$的双变量函数
\end{definition}

\begin{theorem}[叠加定理]
    \[r[m]=\sum_{n=-\infty}^{\infty}k[m,n]e[n]\]
    其中$k[m,n]$为系统的核函数。
\end{theorem}
\begin{proof}
    离散信号也可以视为它与单位脉冲信号的卷积：
    \[e[m]=(e*\delta)[m]=\sum_{n=-\infty}^{\infty}e[n]\delta[m-n]\]
    根据系统的线性性，有
    \begin{align*}
        r[m]&=H(e[m])\\
        &=H_m\left(\sum_{n=-\infty}^{\infty}e[n]\delta[m-n]\right)\\
        &=\sum_{n=-\infty}^{\infty}e[n]H_m(\delta[m-n])
    \end{align*}
    令$k[m,n]=H_m(\delta[m-n])$即得。
\end{proof}

另一方面，由加权求和描述的系统显然也是线性的，因此我们可以得到：
\begin{proposition}[离散线性系统的等价定义]
    一个离散系统是线性系统当且仅当响应是激励的加权求和。
\end{proposition}

类似连续系统，只要系统有线性性，就可以将系统原有能量引起的\textbf{零输入响应}$r_{zi}[n]$与单纯由激励引起的\textbf{零状态响应}$r_{zs}[n]$分开考虑。进一步，如果系统还具有时不变性，则可定义：
\begin{definition}[单位脉冲响应与单位阶跃响应]
    当$e[n]=\delta[n]$时，系统的响应称为\textbf{单位脉冲响应}，记为$h[n]$；当$e[n]=u[n]$时，系统的响应称为\textbf{单位阶跃响应}，记为$s[n]$。
\end{definition}

这样，核函数可以用单位脉冲响应来表示，线性时不变系统的响应可以用卷积和来表示：
\[k[m,n]=H_m[\delta[m-n]]=H(\tau_n \delta[m])=\tau_n H(\delta[m])=h[m-n]\]
\[r[n]=\sum_{m=-\infty}^{\infty}h[n-m]e[m]=(h*e)[n]\]

可见对于离散系统，“线性时不变系统”与“卷积和系统”仍然是等价的概念。

\subsection{离散系统的系统函数}\label{sub:discrete_system_function}

在上一小节中，我们得到$r[n]=(h*e)[n]$，对等式两边做z变换，得到
\[R(z)=H(z)E(z)\]
其中$H(z)=\mathcal{Z} h[n]$记为离散系统的系统函数，它也是响应与激励的z变换之比：
\begin{definition}[离散系统的系统函数]
    设离散系统的激励为$e[n]$，响应为$r[n]$，它们的z变换分别为$E(z),R(z)$，则系统函数定义为
    \[H(z)=\frac{R(z)}{E(z)}\]
\end{definition}

离散系统与连续系统的系统函数分析方法是类似的：一般来说，我们考虑的离散系统的系统函数为有理函数，可以绘制系统函数的零极点图，并通过它来分析系统的滤波特性、稳定性、因果性等时域特性；此外，\textbf{系统流图与梅森增益公式对于离散系统仍然适用}，系统的串联（级联）、并联与反馈的定义与连续系统完全一致，\ref{sub:mason_gain_formula}中介绍的信号流图的标准型也仍然适用于离散系统，因此这部分内容不再单独讨论。

与连续系统不同的是，连续系统的分析通常使用单边拉普拉斯变换，其收敛域是复平面的右边$\text{Re}\ s>\sigma_c$；而离散系统的分析通常使用双边z变换，其收敛域是环状区域$r<|z|<R$。这使得离散系统的稳定性的z域判则不同于连续系统。此外，双边z变换并不要求单位脉冲响应$h[n]$是因果的，因此因果性也具有z域判则。本小节就来探讨离散系统的系统函数与因果性、稳定性之间的关系。对于离散系统的滤波特性，将在下一章介绍离散系统的频率响应特性后进行讨论。

\begin{proposition}[因果性的时域判则]
    一个离散系统是因果系统当且仅当其单位脉冲响应为因果信号，即$h[n]=0,n<0$。
\end{proposition}
\begin{proof}
    类比连续系统，容易理解因果性与上述条件的等价性，首先来验证卷积和也具有类似于卷积的“支集相加”性质：设离散信号$x[n],y[n]$分别在区间$[n_1,n_2]$与$[m_1,m_2]$内非零，而在其他位置均为0，则
    \begin{align*}
        (x*y)[n]=\sum_{k=-\infty}^{\infty}x[k]y[n-k]
    \end{align*}
    \begin{align*}
        (x*y)[n]\neq 0\iff k\in[n_1,n_2],n-k\in[m_1,m_2]\iff n\in[n_1+m_1,n_2+m_2]
    \end{align*}
    
    对于离散系统，如果它是因果的，那么对于任意激励信号$e[n]$，响应$r[n]$都满足$r[n]=0,n<0$，因此单位脉冲响应也满足$h[n]=0,n<0$；反之，如果单位脉冲响应满足$h[n]=0,n<0$，则对于任意激励信号$e[n]$，响应$r[n]=(h*e)[n]$在$n<0$时为0，因此系统是因果的。
\end{proof}

\begin{proposition}[因果性的z域判则]
    一个离散系统是因果系统当且仅当其系统函数的收敛域包含$\infty$。
\end{proposition}

这一判则是显然的，因为因果信号的z变换是只有负幂次项，收敛域包含无穷远点自然也就与单位脉冲响应的因果性等价。

\begin{proposition}[稳定性的时域判则]
    一个离散系统是稳定系统当且仅当$h\in l^1(\mathbb{Z})$。
\end{proposition}
\begin{proof}
    一方面，对于有界激励$e\in l^{\infty}(\mathbb{Z})$，如果单位脉冲响应满足$h[n]\in l^1(\mathbb{Z})$，则
    \begin{align*}
        |r[n]|&=|(h*e)[n]|\\
        &=\left|\sum_{m=-\infty}^{\infty}h[n-m]e[m]\right|\\
        &\leq \sum_{m=-\infty}^{\infty}|h[n-m]||e[m]|\\
        &\leq \|e\|_{\infty}\sum_{m=-\infty}^{\infty}|h[n-m]|\\
        &=\|e\|_{\infty}\|h\|_1<\infty
    \end{align*}
    即$r\in l^{\infty}(\mathbb{Z})$，系统是稳定的。
    
    另一方面，设$h$不是绝对可和的，取有界激励信号$e[n]=\text{sgn}(h[-n])$，则
    \begin{align*}
        |r[0]|&=|(h*e)[0]|\\
        &=\left|\sum_{m=-\infty}^{\infty}h[m]e[-m]\right|\\
        &=\sum_{m=-\infty}^{\infty}|h[m]|\\
        &=\|h\|_1=\infty
    \end{align*}
    $e\in l^{\infty}(\mathbb{Z})$，但$r\notin l^{\infty}(\mathbb{Z})$，系统不稳定。
\end{proof}

对稳定性的时域判据做z变换，立即得到：

\begin{proposition}[稳定性的z域判则]
    一个离散系统是稳定系统当且仅当其系统函数的收敛域包含单位圆。
\end{proposition}

综上，对于离散线性时不变系统，因果且稳定就等价于其系统函数的收敛域包含以下区域：
\[\left\{\left.z\in\hat{\mathbb{C}}\right||z|\geq 1\right\}\]
